% Gauss, Werke X/1 (1917), native137 lower–141, printed130–134.
\NumberedItem{2}
\begin{center}\large\textsc{Gauss} an \textsc{Schumacher}.\quad Göttingen, 2. April 1833.\end{center}
\HistoricCitation{Briefwechsel zwischen \textsc{Gauss} und \textsc{Schumacher}. II. Altona 1860, S. 328.}
\noindent\hspace*{\parindent}Recht sehr danke ich Ihnen, theuerster Freund, für die gütige Mittheilung
des lineirten Papiers. Ich finde es sehr brauchbar besonders für Zeichnungen
jeder Art, denen rechtwinklige Coordinaten zum Grunde liegen, insofern
nicht die grösste Genauigkeit gefordert wird, und das ist bei meinen
Zeichnungen ohnehin nie der Fall, da ich sie nie mache, um etwas durch
Abmessung definitiv daraus abzuleiten; also für Zeichnung terrestrischer Puncte,
allenfalls auch selbst kleiner Sternkarten, Versinnlichung des Ganges des Barometers,
Variationen der Magnetnadel etc. etc., nicht weniger auch zu Zeichnungen,
die sich auf rein mathematische Sachen beziehen, wie z. B. in Rücksicht
auf die imaginären Wurzeln der Gleichungen, und gerade mit Gegenständen
der letztern Art habe ich mich in der letzten Zeit viel beschäftigt.
\clearpage
\SourceLeaf{138}\ArithmeticHead{131}{R}{GAUSS AN SCHUMACHER.}
\noindent Da Sie jetzt wohl meine Anzeige von \textsc{Fourier}[*)] gelesen haben, so interessirt
Sie vielleicht die Bemerkung, dass ich absichtlich S. 324[**)] »sondern
so lange als .\,.\,.\,.\,.\,. zweifelhaft bleiben muss« hier in gemessenen Worten
mich ausgedrückt habe, aber nicht weil ich \textls[120]{selbst} über Dasein oder Nichtdasein
eines solchen Zusammenhangs ungewiss geblieben, sondern weil in den
G. G. A. nicht der Ort war mich darüber auf bestimmtere Weise zu erklären.
Ich glaube auf das klarste nachweisen zu können, dass ein solcher Zusammenhang
\textls[120]{nicht} existirt, allein dies wird erst geschehen können, wenn ich einmal
Gelegenheit nehme meine Untersuchungen über die Wurzeln der Gleichungen
ausgearbeitet bekannt zu machen. Sie wissen, dass ich langsam schreibe,
allein das kommt hauptsächlich daher, weil ich mir nie anders gefallen kann,
als wenn in kleinem Raum möglichst viel ist, und kurz zu schreiben viel mehr
Zeit kostet als lang. Wollte ich jene Untersuchungen, die, wenn ich sie
einmal entwickele, nur eine mässige Zahl Bogen betragen dürfen, mit der
Breite, wie \textsc{Fouriers} Buch geschrieben ist, vortragen, so würde ich vielleicht
nur $\tfrac14$ so viel Zeit und mehrere grosse Quartbände gebrauchen.
.\,.\,.\,.\,.\,.
\NumberedItem{3}
\begin{center}\large\textsc{Gauss} an \textsc{Schumacher}.\quad Göttingen, 1836. Junius 20.\end{center}
\HistoricCitation{Briefwechsel zwischen \textsc{Gauss} und \textsc{Schumacher}. III. Altona 1861, S. 68--69.}
\noindent .\,.\,.\,.\,.\,.\,.\,.\quad Bei Gelegenheit der Vorlesung, die ich halte, bin ich veranlasst
in diesen Tagen auf die Theorie der Gleichungen zurückzukommen, der ich
jetzt einen ganz neuen Gesichtspunkt abgewonnen habe, in Folge von welchem
es mir nicht unwahrscheinlich ist, dass sich \textls[120]{dennoch} eine Brücke zwischen
\textsc{Fouriers} critischen Punkten und bestimmten Paaren von imaginären Wurzeln
schlagen lässt. Vor etwa 3 Jahren habe ich in einem Briefe an Sie die entgegengesetzte
Meinung ausgesprochen, aber freilich hatte ich damals nicht
alle auf den Gegenstand bezüglichen Details entwickelt, was, wenn man einen
\par\smallskip\noindent\rule{.25\linewidth}{.4pt}\par\smallskip
{\small\noindent\hspace*{\parindent}[*) Werke III, S. 119.]\par
\noindent\hspace*{\parindent}[**) Werke III, S. 121, Zeile 10--12.]\par}
\par\vfill\hfill17*\par
\clearpage
\SourceLeaf{139}\ArithmeticHead{132}{L}{ALGEBRA.\quad NACHLASS.}
\noindent negativen Satz zu beweisen hat, nothwendig ist um ganz gewiss zu sein. Die
allgemeinen Apperçus, die in 999 Fällen gut gehen, können das 1000te mal am
Ende in einen cul de sac führen. Berichtigen Sie diesmal meine damaligen
vielleicht zu positiven Behauptungen, aber bemerken Sie zugleich, dass ich
nicht ohne Ursach pauca sed matura zu meinem Wahlspruch für alles zu
\textls[120]{veröffentlichende} gemacht habe. Die allgemeinen Apperçus sind die Geburten
Einer Stunde, aber um daraus etwas gereiftes zu machen, ist oft lange,
oft jahrelange grosse Detailarbeit nöthig, von der man voraussieht, dass man
sie gewiss durchführen kann, wenn man sich dazu gibt, obwohl auch dann
immer noch manche ähnliche Geburten zweiten und dritten Ranges, die schon
auf Ordre kommen müssen, nöthig sind. Procreare iucundum, at parturire
molestum. Was übrigens meinen gegenwärtigen Gesichtspunkt betrifft, so
würden auch erst sehr zeitraubende Details, um alle Verzweigungen zu verfolgen,
nöthig sein, und ich glaube, nicht, dass ich, indem so vielerlei unter
\par\noindent\makebox[\linewidth][s]{. . . . . . . . . . . . . . . . . . . . . . . . . . .}\par
\begin{center}\small[Der Schluß des Briefes fehlt.]\end{center}
\NumberedItem{4}
\begin{center}\large\textsc{Gauss} an \textsc{Schumacher}.\quad Göttingen, den 24. Juni 1836.\end{center}
\HistoricCitation{Briefwechsel zwischen \textsc{Gauss} und \textsc{Schumacher}. III. Altona 1861. S. 72.}
\noindent .\,.\,.\,.\,.\,.\quad Ich habe in den letzten Tagen meine Ideen über die Gleichungen
weiterverfolgt, aber das Resultat ist eher das entgegengesetzte, und führt vielmehr
zur Bestärkung in meiner früheren Ansicht, dass es einen in der Natur
der Sache liegenden allgemeinen willkürfreien Zusammenhang zwischen den
einzelnen critischen Punkten und den einzelnen Paaren von imaginären Wurzeln
gar nicht gibt. Es scheint mir selbst, dass ich in meinem letzten Briefe dem
Ausdruck meiner früheren Ansicht in einem älteren Briefe Unrecht gethan
habe, indem mir jetzt (nach langer Unterbrechung) jene frühere Ansicht \textls[120]{nicht}
\textls[120]{gleich} in derselben Frische gegenwärtig war, in der ich sie damals aufgefasst
hatte. Soviel bleibt aber immer gewiss, dass bei solchen negativen Sätzen [für]
die Verwandlung der subjectiven Überzeugung in eine objective (für andere)
\clearpage
\SourceLeaf{140}\ArithmeticHead{133}{R}{GAUSS AN SCHUMACHER.}
\noindent eine höchst abschreckende Detailarbeit erforderlich wäre. Man würde, um
die Verschiedenheit der Fälle wirklich anschaulich zu machen, eine grosse
Menge von Gleichungen in concreto durch Curven versinnlichen müssen; jede
Curve müsste durch Punkte gezeichnet werden, und die Bestimmung eines
einzigen Punkts erfordert schon langwierige Rechnungen. Sie sehen es wohl
der Fig. 4 bei meiner ersten Schrift von 1799 nicht an, wie viel Arbeit die
richtige Zeichnung dieser Curve erfordert hat, und doch ist diess vergleichungsweise
nur ein sehr einfacher Fall gegen viele, die hier betrachtet werden
müssten. .\,.\,.\,.\,.\,.
\par\bigskip\begin{center}\textbf{BEMERKUNGEN.}\end{center}
{\small
\noindent\hspace*{\parindent}\textbf{I.} Von der Richtigkeit des von \textsc{Gauss} in der Aufzeichnung [1.] ausgesprochenen Satzes kann man sich
nach einer mündlichen Mitteilung von F. \textsc{Klein} wie folgt überzeugen. Setzt man $f(x)=w$ und bezeichnet
die Wurzeln der Gleichung $f'(x)=0$ — die der Einfachheit wegen alle von einander verschieden sein mögen —
durch $\alpha_1,\alpha_2,\ldots,\alpha_{n-1}$, so sind die Punkte $w=f(\alpha_1),f(\alpha_2),\ldots,f(\alpha_{n-1})$ die im Endlichen gelegenen einfachen
Verzweigungspunkte der algebraischen Funktion $x$ von $w$. Denkt man sich die zu dieser Funktion gehörige
$n$-blättrige \textsc{Riemann}sche Fläche in der Weise konstruiert, daß man durch die Verzweigungspunkte $f(\alpha_1)$,
$f(\alpha_2),\ldots,f(\alpha_{n-1})$ zur reellen $w$-Achse parallel laufende Verzweigungsschnitte legt und etwa längs des von
$f(\alpha_\varkappa)$ ausgehenden Schnittes die Blätter $1$ und $\varkappa+1$ aneinander heftet, so liefert die Abbildung der Blätter
dieser Fläche auf die $x$-Ebene die durch die Bilder der Verzweigungsschnitte begrenzten $n$ Flächenräume,
von denen in dem \textsc{Gauss}schen Satze die Rede ist. In der Handschrift geht diesem Satze unmittelbar voran
der Werke III, S. 112 und Werke VIII, S. 32 abgedruckte Lehrsatz mit einem Beispiel. Beide Sätze gehören
zusammen und fallen, wie auch die Eingangsworte der Aufzeichnung [1.] zeigen, unter den »ganz
neuen Gesichtspunkt«, den \textsc{Gauss} in der Briefstelle [3.] erwähnt.
\par\hfill\textsc{Schlesinger}.\par\medskip
\noindent\hspace*{\parindent}\textbf{II.} Die in den Briefstellen [2.]—[4.] an \textsc{Schumacher} und auch schon früher in dem Briefe an \textsc{Drobisch}
(S. 106, 107) gemachten Bemerkungen, die an \textsc{Fouriers} \textit{Analyse des équations déterminées} (Paris,
1831, deutsch in \textsc{Ostwalds} Klassikern Nr. 127) anknüpfen, verlangen ein Eingehen auf \textsc{Gauss'} Anzeige dieses
Werkes (Werke III, S. 119), da einige Angaben dieser Anzeige einer Richtigstellung bedürfen. Die Aussage
Werke III, S. 121, Zeile 4: »In der That ist es zwar wahr, dass die Gleichung $X=0$ \textls[120]{zusammengezählt}
genau soviele Paare imaginärer Wurzeln enthält, als solche Ausfälle oder kritische Stellen vorkommen«,
trifft bei der a. a. O. S. 120, Zeile 5 bis 1 v. u. gegebenen Definition der kritischen Stellen nicht
zu, wie schon das Beispiel der Gleichung $(x-b)^n-c=0$ zeigt, wo $n$ eine gerade Zahl $\geqq4$, $b$ beliebig
reell, $c$ positiv sein soll. Trotz ihrer $n-2$ verschiedenen imaginären Wurzeln hätte die Gleichung keine
kritische Stelle im Sinne von \textsc{Gauss}. Wir definieren (vergl. Archiv der Mathem. und Physik (3) \textbf{16}, 1910,
S. 157) eine kritische Stelle auf folgende Weise: Es bedeuten $f(x)=0$ eine Gleichung $n^{\mathrm{ten}}$ Grades mit
reellen Koeffizienten, $f'(x),f''(x),\ldots$ die Derivierten von $f(x)$ und $\alpha$ eine reelle Größe; in der Reihe $f(\alpha)$,
$f'(\alpha),\ldots,f^{(n)}(\alpha)$ sei $e$ die Anzahl der verschwindenden intermediären Funktionen *), $c$ und $d$ die Anzahl
\par\smallskip\noindent\rule{.25\linewidth}{.4pt}\par\smallskip
\noindent\hspace*{\parindent}*) Sollte also etwa $f(\alpha)=f'(\alpha)=\cdots=f^{(l-1)}(\alpha)=0$ sein, so sind diese ersten $l$ verschwindenden
Größen bei der Bestimmung von $e$ nicht mitzuzählen.
\par}
\clearpage
\SourceLeaf{141}\ArithmeticHead{134}{L}{ALGEBRA.\quad NACHLASS.}
{\small\noindent der durch eine ungerade Zahl fehlender Glieder unterbrochenen Zeichenwechsel bezw. Zeichenfolgen und
$\sigma=e-c+d$. Wenn $\sigma\geqq2$ ist, so heiße $\alpha$ eine kritische Stelle von $f(x)=0$ *). Bei dieser Definition
gilt der folgende \textls[120]{präzisierte} \textsc{Fourier}sche Satz: Es seien $a$ und $a'$ zwei beliebige reelle Zahlen $(a<a')$;
bedeutet $Z_a$ bezw. $Z_{a'}$ die Anzahl der Vorzeichenwechsel in den Reihen
\[
\begin{array}{r@{\qquad}l}
(1)&f(a),\ f'(a),\ f''(a),\ \ldots,\ f^{(n)}(a),\\[.5ex]
(2)&f(a'),\ f'(a'),\ f''(a'),\ \ldots,\ f^{(n)}(a')
\end{array}
\]
und ist $2\lambda$ die Summe der sämtlichen ihrer Definition nach niemals ungeraden Zahlen $\sigma$, die allen zwischen
$a$ und $a'$ gelegenen kritischen Stellen von $f(x)=0$ entsprechen, so ist stets
\[
\begin{array}{r@{\qquad}l}(3)&Z_a-Z_{a'}-2\lambda\geqq0\end{array}
\]
und die durch die linke Seite von (3) definierte Zahl gibt die \textls[120]{genaue} Anzahl der zwischen den Grenzen
$a$ und $a'$ gelegenen Wurzeln von $f(x)=0$, jede Wurzel nach ihrer Vielfachheit gezählt. Für den Fall,
daß die Grenzen $a$ und $a'$ kritische Stellen oder Wurzeln von $f(x)=0$ sind, ist für die Bildung von
$2\lambda$ und die Zählung der Gleichungswurzeln $a$ auszuschließen, $a'$ einzuschließen **). Für $a=-\infty$,
$a'=+\infty$ erhält man den folgenden Satz, der an die Stelle der \textsc{Gauss}schen Angabe zu treten hat: Die
Anzahl aller imaginären Wurzeln von $f(x)=0$ ist gleich der Summe aller Zahlen $\sigma$, die den sämtlichen
zwischen $-\infty$ und $+\infty$ gelegenen kritischen Stellen von $f(x)=0$ entsprechen. — Aus der Ungleichung
(3) folgt, daß, wenn zwischen $a$ und $a'$ eine kritische Stelle liegt, für die $\sigma>2$ ist, $Z_a-Z_{a'}>2$ sein muß.
Hiernach trifft die Aussage von \textsc{Gauss} (Werke III, S. 120, Zeile 11 v. u.), wonach »sämtliche einzelne
Differenzen $g'-g,\ g''-g',\ g'''-g''$ u. s. w. nur entweder $=1$ oder $=2$ werden«, wenn man die Grenzen
eng genug wählt, nicht zu. Daß man $g'-g$ nichts stets $\leqq2$ machen kann, war \textsc{Fourier} bekannt, wie
aus seiner »Regel des doppelten Vorzeichens« (a. a. O. S. 103) hervorgeht, wo er auch das Beispiel einer
Gleichung mit 14 imaginären Wurzeln erwähnt. Die oben gegebene Definition der kritischen Stellen ist
ihrem Wesen nach auch nur eine Präzisierung der genannten \textsc{Fourier}schen Regel.

Daß sich nicht »eine Brücke zwischen \textsc{Fouriers} critischen Punkten und bestimmten Paaren von imaginären
Wurzeln schlagen lässt« (vergl. die Briefstelle [3.]), folgt aus dem Beispiel $(x-b)^n+c=0$ ($n$ gerade,
$b$ reell, $c>0$), mit der einzigen kritischen Stelle $b$ und $n$ verschiedenen imaginären Wurzeln. Offen bleibt
aber noch die folgende Frage: Es sei für $f(x)=0$ die Anzahl $m$ der kritischen Stellen nicht größer als $\dfrac n2$
und für jede $\sigma=2$ ***); gibt es dann »einen willkürfreien Zusammenhang« (vergl. die Briefstelle [4.]) zwischen
den einzelnen dieser $m$ kritischen Stellen und $m$ Paaren von imaginären Wurzeln der Gleichung $f(x)=0$?
\par\hfill\textsc{Alfred Loewy}.\par
\smallskip\noindent\rule{.25\linewidth}{.4pt}\par\smallskip
\noindent\hspace*{\parindent}*) Für das Beispiel ist $b$ eine kritische Stelle mit $e=n-1$, $c=1$, $d=0$, also $\sigma=n-2$.
\par\noindent\hspace*{\parindent}**) Der von \textsc{Gauss}, Werke III, S. 120, Zeile 8, bei Seite gesetzte »ausnahmliche Fall« schließt aus,
daß eine der Größen aus den Reihen (1), (2) gleich Null wird. Es handelt sich also um eine etwas stärkere
Beschränkung, als wenn man fordert, daß $a$ und $a'$ weder Gleichungswurzeln noch kritische Stellen sein sollen.
\par\noindent\hspace*{\parindent}***) Schon der Fall $\sigma=2$ kommt nicht nur entsprechend der \textsc{Gauss}schen Anzeige, die nur $d=1$
zuläßt, zustande, sondern in allgemeinster Weise für ein beliebiges ganzzahliges $e$, wenn $c=e-1$, $d=1$
oder $c=e-2$, $d=0$ ist. Es folgt dies aus $\sigma=2$ und aus der gemäß Definition der Zahlen $c,d,e$
stets giltigen Relation $e\geqq c+d$.
\par}
\ClosingRule
